% ==================================================================
% id:       prob-014
% topic:    מבוא להסתברות
% subtopic: עקרון ההכלה וההפרדה
% solution: full
% ==================================================================
חברת פרסום ערכה סקר כדי לבדוק כמה אנשים קוראים את העיתונים „ידיעות אחרונות” ו„מעריב”. לפי הסקר, 55\% קוראים את „ידיעות אחרונות”, 50\% קוראים את „מעריב”, ו-25\% קוראים את שניהם. מה אחוז האנשים ש:
\begin{enumerate}
\item קוראים את „ידיעות אחרונות” או את „מעריב”?
\begin{solbox}
נסמן $Y$ = „קורא ידיעות”, $M$ = „קורא מעריב”. לפי הכלה והפרדה: $\PP(Y\cup M)=0.55+0.5-0.25=0.8$, כלומר $\ans{80\%}$.
\end{solbox}
\item לא קוראים את „ידיעות אחרונות”?
\begin{solbox}
$\PP(Y^C)=1-0.55=0.45$, כלומר $\ans{45\%}$.
\end{solbox}
\item לא קוראים אף אחד משני העיתונים?
\begin{solbox}
$\PP\big((Y\cup M)^C\big)=1-0.8=0.2$, כלומר $\ans{20\%}$.
\end{solbox}
\item קוראים רק את „מעריב”?
\begin{solbox}
$\PP(M\cap Y^C)=\PP(M)-\PP(M\cap Y)=0.5-0.25=0.25$, כלומר $\ans{25\%}$.
\end{solbox}
\item קוראים בדיוק עיתון אחד?
\begin{solbox}
$\PP(Y\cup M)-\PP(Y\cap M)=0.8-0.25=0.55$, כלומר $\ans{55\%}$.
\end{solbox}
\item לא קוראים את שני העיתונים יחד (כלומר, יש לפחות עיתון אחד מהשניים שהם לא קוראים)?
\begin{solbox}
זה המשלים של „קוראים את שניהם”: $1-0.25=0.75$, כלומר $\ans{75\%}$.
\end{solbox}
\end{enumerate}
